\documentclass{amsart}
% For dealing with references we use the comment environment
\usepackage{verbatim}
\newenvironment{reference}{\comment}{\endcomment}
%\newenvironment{reference}{}{}
\newenvironment{slogan}{\comment}{\endcomment}
\newenvironment{history}{\comment}{\endcomment}

% For commutative diagrams we use Xy-pic
\usepackage[all]{xy}

% We use 2cell for 2-commutative diagrams.
\xyoption{2cell}
\UseAllTwocells

% We use multicol for the list of chapters between chapters
\usepackage{multicol}

% This is generally recommended for better output
\usepackage{lmodern}
\usepackage[T1]{fontenc}

% For cross-file-references

% Package for hypertext links:
\usepackage{hyperref}

% For any local file, say "hello.tex" you want to link to please

% Theorem environments.
%
\theoremstyle{plain}
\newtheorem{theorem}[subsection]{Theorem}
\newtheorem{proposition}[subsection]{Proposition}
\newtheorem{lemma}[subsection]{Lemma}

\theoremstyle{definition}
\newtheorem{definition}[subsection]{Definition}
\newtheorem{example}[subsection]{Example}
\newtheorem{exercise}[subsection]{Exercise}
\newtheorem{situation}[subsection]{Situation}

\theoremstyle{remark}
\newtheorem{remark}[subsection]{Remark}
\newtheorem{remarks}[subsection]{Remarks}

\numberwithin{equation}{subsection}

% Macros
%
\def\lim{\mathop{\mathrm{lim}}\nolimits}
\def\colim{\mathop{\mathrm{colim}}\nolimits}
\def\Spec{\mathop{\mathrm{Spec}}}
\def\Hom{\mathop{\mathrm{Hom}}\nolimits}
\def\Ext{\mathop{\mathrm{Ext}}\nolimits}
\def\SheafHom{\mathop{\mathcal{H}\!\mathit{om}}\nolimits}
\def\SheafExt{\mathop{\mathcal{E}\!\mathit{xt}}\nolimits}
\def\Sch{\mathit{Sch}}
\def\Mor{\mathop{\mathrm{Mor}}\nolimits}
\def\Ob{\mathop{\mathrm{Ob}}\nolimits}
\def\Sh{\mathop{\mathit{Sh}}\nolimits}
\def\NL{\mathop{N\!L}\nolimits}
\def\CH{\mathop{\mathrm{CH}}\nolimits}
\def\proetale{{pro\text{-}\acute{e}tale}}
\def\etale{{\acute{e}tale}}
\def\QCoh{\mathit{QCoh}}
\def\Ker{\mathop{\mathrm{Ker}}}
\def\Im{\mathop{\mathrm{Im}}}
\def\Coker{\mathop{\mathrm{Coker}}}
\def\Coim{\mathop{\mathrm{Coim}}}

% Boxtimes
%
\DeclareMathSymbol{\boxtimes}{\mathbin}{AMSa}{"02}

%
% Macros for moduli stacks/spaces
%
\def\QCohstack{\mathcal{QC}\!\mathit{oh}}
\def\Cohstack{\mathcal{C}\!\mathit{oh}}
\def\Spacesstack{\mathcal{S}\!\mathit{paces}}
\def\Quotfunctor{\mathrm{Quot}}
\def\Hilbfunctor{\mathrm{Hilb}}
\def\Curvesstack{\mathcal{C}\!\mathit{urves}}
\def\Polarizedstack{\mathcal{P}\!\mathit{olarized}}
\def\Complexesstack{\mathcal{C}\!\mathit{omplexes}}
% \Pic is the operator that assigns to X its picard group, usage \Pic(X)
% \Picardstack_{X/B} denotes the Picard stack of X over B
% \Picardfunctor_{X/B} denotes the Picard functor of X over B
\def\Pic{\mathop{\mathrm{Pic}}\nolimits}
\def\Picardstack{\mathcal{P}\!\mathit{ic}}
\def\Picardfunctor{\mathrm{Pic}}
\def\Deformationcategory{\mathcal{D}\!\mathit{ef}}

\setlength{\emergencystretch}{3em}
\begin{document}
\title[Stacks Project, Tag 032P]{1271 --- Tate's Japanese ascent for x-adically complete normal Noetherian domains with Japanese domain quotient}
\maketitle
\noindent Copyright (C) 2005--2025 Johan de Jong. Stacks Project authors. Source: \href{https://stacks.math.columbia.edu/tag/032P}{Tag 032P}.

\noindent Editorial difficulty note (not author text). Difficulty H2: The author reduces normalization finiteness to purely inseparable fraction-field extensions, adjoins a root of the completion parameter, identifies a principal valuation prime, and lifts finite residue normalization through separated adic modules. This coupled inseparability, one-dimensional valuation and completion argument exceeds ordinary qualification commutative algebra..

\noindent Editorial provenance note (not author text). History: excerpt from revision \texttt{a0444\allowbreak 6e57e\allowbreak c1fbc\allowbreak 252a8\allowbreak 71afc\allowbreak ec775\allowbreak 2fb28\allowbreak 07b14}, algebra.tex, lines 45534--45590. Reference presentation was adapted; original-author context and core support blocks were retained or added. The mathematical wording is unchanged. Licensed under GNU FDL 1.2 or later; no invariant sections or cover texts. See ../licenses/COPYING. Context blocks reproduce author definitions and supporting results; they are not additional counted problems.

\begin{definition}
\label{definition-N}
\begin{reference}
\href{https://stacks.math.columbia.edu/bibliography}{Bibliography: EGA (Chapter 0, Definition 23.1.1)}
\end{reference}
Let $R$ be a domain with field of fractions $K$.
\begin{enumerate}
\item We say $R$ is {\it N-1} if the integral closure of $R$ in $K$
is a finite $R$-module.
\item We say $R$ is {\it N-2} or {\it Japanese} if for any finite
extension $L/K$ of fields the integral closure of $R$ in $L$
is finite over $R$.
\end{enumerate}
\end{definition}

\begin{definition}
\label{fields-definition-trace-pairing}
Let $L/K$ be a finite extension of fields. The {\it trace pairing}
for $L/K$ is the symmetric $K$-bilinear form
$$
Q_{L/K} : L \times L \longrightarrow K,\quad
(\alpha, \beta) \longmapsto \text{Trace}_{L/K}(\alpha\beta)
$$
\end{definition}

\begin{lemma}
\label{lemma-Noetherian-normal-domain-finite-separable-extension}
Let $R$ be a Noetherian normal domain with fraction field $K$.
Let $L/K$ be a finite separable field extension.
Then the integral closure of $R$ in $L$ is finite over $R$.
\end{lemma}

\begin{proof}
Consider the trace pairing
(Fields, Definition \ref{fields-definition-trace-pairing})
$$
L \times L \longrightarrow K,
\quad (x, y) \longmapsto \langle x, y\rangle := \text{Trace}_{L/K}(xy).
$$
Since $L/K$ is separable this is nondegenerate
(Fields, Lemma \href{https://stacks.math.columbia.edu/tag/0BIL}{lemma separable trace pairing (Tag 0BIL)}).
Moreover, if $x \in L$ is integral over $R$, then
$\text{Trace}_{L/K}(x)$ is in $R$. This is true because the
minimal polynomial of $x$ over $K$ has coefficients in $R$
(Lemma \href{https://stacks.math.columbia.edu/tag/00H7}{lemma minimal polynomial normal domain (Tag 00H7)})
and because $\text{Trace}_{L/K}(x)$ is an
integer multiple of one of these coefficients
(Fields, Lemma \href{https://stacks.math.columbia.edu/tag/0BIH}{lemma trace and norm from minimal polynomial (Tag 0BIH)}).
Pick $x_1, \ldots, x_n \in L$ which are integral over $R$
and which form a $K$-basis of $L$. Then the integral closure
$S \subset L$ is contained in the $R$-module
$$
M = \{y \in L \mid \langle x_i, y\rangle \in R, \ i = 1, \ldots, n\}
$$
By linear algebra we see that $M \cong R^{\oplus n}$ as an $R$-module.
Hence $S \subset R^{\oplus n}$ is a finitely generated $R$-module
as $R$ is Noetherian.
\end{proof}


\begin{lemma}
\label{lemma-domain-char-p-N-1-2}
Let $R$ be a Noetherian domain with fraction field $K$ of
characteristic $p > 0$. Then $R$ is N-2 if and only if
for every finite purely inseparable extension $L/K$ the integral
closure of $R$ in $L$ is finite over $R$.
\end{lemma}

\begin{proof}
Assume the integral closure of $R$ in every finite purely inseparable
field extension of $K$ is finite.
Let $L/K$ be any finite extension. We have to show the
integral closure of $R$ in $L$ is finite over $R$.
Choose a finite normal field extension $M/K$
containing $L$. As $R$ is Noetherian it suffices to show that
the integral closure of $R$ in $M$ is finite over $R$.
By Fields, Lemma \href{https://stacks.math.columbia.edu/tag/030M}{lemma normal case (Tag 030M)}
there exists a subextension $M/M_{insep}/K$
such that $M_{insep}/K$ is purely inseparable, and $M/M_{insep}$
is separable. By assumption the integral closure $R'$ of $R$ in
$M_{insep}$ is finite over $R$. By
Lemma \ref{lemma-Noetherian-normal-domain-finite-separable-extension}
the integral
closure $R''$ of $R'$ in $M$ is finite over $R'$. Then $R''$ is finite
over $R$ by Lemma \href{https://stacks.math.columbia.edu/tag/00GL}{lemma finite transitive (Tag 00GL)}.
Since $R''$ is also the integral closure
of $R$ in $M$ (see Lemma \href{https://stacks.math.columbia.edu/tag/0308}{lemma integral closure transitive (Tag 0308)}) we win.
\end{proof}


\begin{lemma}
\label{lemma-finite-length}
Let $R$ be a domain with fraction field $K$.
Let $M$ be an $R$-submodule of $K^{\oplus r}$.
Assume $R$ is local Noetherian of dimension $1$.
For any nonzero $x \in R$ we have $\text{length}_R(R/xR) < \infty$
and
$$
\text{length}_R(M/xM) \leq r \cdot \text{length}_R(R/xR).
$$
\end{lemma}

\begin{proof}
If $x$ is a unit then the result is true. Hence we may assume
$x \in \mathfrak m$ the maximal ideal of $R$. Since $x$ is not
zero and $R$ is a domain we have $\dim(R/xR) = 0$, and hence
$R/xR$ has finite length. Consider $M \subset K^{\oplus r}$ as
in the lemma. We may assume that the elements of $M$ generate
$K^{\oplus r}$ as a $K$-vector space after replacing $K^{\oplus r}$
by a smaller subspace if necessary.

\medskip\noindent
Suppose first that $M$ is a finite $R$-module. In that case we can clear
denominators and assume $M \subset R^{\oplus r}$. Since
$M$ generates $K^{\oplus r}$ as a vectors space we see that
$R^{\oplus r}/M$ has finite length. In particular there exists
an integer $c \geq 0$ such that $x^cR^{\oplus r} \subset M$.
Note that $M \supset xM \supset x^2M \supset \ldots$ is a sequence of
modules with successive quotients each isomorphic to $M/xM$. Hence
we see that
$$
n \text{length}_R(M/xM) = \text{length}_R(M/x^nM).
$$
The same argument for $M = R^{\oplus r}$ shows that
$$
n \text{length}_R(R^{\oplus r}/xR^{\oplus r}) =
\text{length}_R(R^{\oplus r}/x^nR^{\oplus r}).
$$
By our choice of $c$ above we see that $x^nM$ is sandwiched between
$x^n R^{\oplus r}$ and $x^{n + c}R^{\oplus r}$. This easily gives that
$$
r(n + c) \text{length}_R(R/xR)
\geq
n \text{length}_R(M/xM)
\geq
r (n - c) \text{length}_R(R/xR)
$$
Hence in the finite case we actually get the result of the lemma with
equality.

\medskip\noindent
Suppose now that $M$ is not finite. Suppose that the length of $M/xM$ is
$\geq k$ for some natural number $k$. Then we can find
$$
0 \subset N_0 \subset N_1 \subset N_2 \subset \ldots N_k \subset M/xM
$$
with $N_i \not = N_{i + 1}$ for $i = 0, \ldots k - 1$.
Choose an element $m_i \in M$ whose congruence class mod $xM$ falls
into $N_i$ but not into $N_{i - 1}$ for $i = 1, \ldots, k$.
Consider the finite $R$-module $M' = Rm_1 + \ldots + Rm_k \subset M$.
Let $N'_i \subset M'/xM'$ be the inverse image of $N_i$.
It is clear that $N'_i \not =N'_{i + 1}$ by our choice of $m_i$.
Hence we see that $\text{length}_R(M'/xM') \geq k$. By the
finite case we conclude $k \leq r\text{length}_R(R/xR)$
as desired.
\end{proof}


\begin{lemma}
\label{lemma-finite-extension-residue-fields-dimension-1}
Let $R \to S$ be a homomorphism of domains inducing an
injection of fraction fields $K \subset L$. If $R$ is Noetherian
local of dimension $1$ and $[L : K] < \infty$ then
\begin{enumerate}
\item each prime ideal $\mathfrak n_i$ of $S$ lying over
the maximal ideal $\mathfrak m$ of $R$ is maximal,
\item there are finitely many of these, and
\item $[\kappa(\mathfrak n_i) : \kappa(\mathfrak m)] < \infty$ for each $i$.
\end{enumerate}
\end{lemma}

\begin{proof}
Pick $x \in \mathfrak m$ nonzero. Apply Lemma \ref{lemma-finite-length}
to the submodule $S \subset L \cong K^{\oplus n}$ where $n = [L : K]$.
Thus the ring $S/xS$ has finite length over $R$. It follows that
$S/\mathfrak m S$ has finite length over $\kappa(\mathfrak m)$.
In other words, $\dim_{\kappa(\mathfrak m)} S/\mathfrak m S$
is finite (Lemma \href{https://stacks.math.columbia.edu/tag/00IY}{lemma dimension is length (Tag 00IY)}). Thus $S/\mathfrak mS$
is Artinian (Lemma \href{https://stacks.math.columbia.edu/tag/00J6}{lemma finite dimensional algebra (Tag 00J6)}). The
structural results on Artinian rings implies parts (1) and (2), see
for example Lemma \href{https://stacks.math.columbia.edu/tag/00JB}{lemma artinian finite length (Tag 00JB)}.
Part (3) is implied by the finiteness established above.
\end{proof}


\begin{lemma}
\label{lemma-finite-length-global}
Let $R$ be a domain with fraction field $K$.
Let $M$ be an $R$-submodule of $K^{\oplus r}$.
Assume $R$ is Noetherian of dimension $1$.
For any nonzero $x \in R$ we have
$\text{length}_R(M/xM) < \infty$.
\end{lemma}

\begin{proof}
Since $R$ has dimension $1$ we see that $x$ is contained in finitely many
primes $\mathfrak m_i$, $i = 1, \ldots, n$, each maximal. Since $R$ is
Noetherian we see that $R/xR$ is Artinian and
$R/xR = \prod_{i = 1, \ldots, n} (R/xR)_{\mathfrak m_i}$ by
Proposition \href{https://stacks.math.columbia.edu/tag/00KJ}{proposition dimension zero ring (Tag 00KJ)} and
Lemma \href{https://stacks.math.columbia.edu/tag/00JB}{lemma artinian finite length (Tag 00JB)}. Hence $M/xM$ similarly
decomposes as the product $M/xM = \prod (M/xM)_{\mathfrak m_i}$
of its localizations at the $\mathfrak m_i$. By Lemma \ref{lemma-finite-length}
applied to $M_{\mathfrak m_i}$ over $R_{\mathfrak m_i}$ we see each
$M_{\mathfrak m_i}/xM_{\mathfrak m_i} = (M/xM)_{\mathfrak m_i}$
has finite length over $R_{\mathfrak m_i}$. Thus $M/xM$
has finite length over $R$ as the above implies $M/xM$
has a finite filtration by $R$-submodules whose successive
quotients are isomorphic to the residue fields $\kappa(\mathfrak m_i)$.
\end{proof}


\begin{lemma}[Krull-Akizuki]
\label{lemma-krull-akizuki}
Let $R$ be a domain with fraction field $K$.
Let $L/K$ be a finite extension of fields.
Assume $R$ is Noetherian and $\dim(R) = 1$.
In this case any ring $A$ with $R \subset A \subset L$ is
Noetherian.
\end{lemma}

\begin{proof}
Let $I \subset A$ be a nonzero ideal. By
Lemma \href{https://stacks.math.columbia.edu/tag/0H7L}{lemma intersection not zero (Tag 0H7L)}
we can find a nonzero element $x \in I \cap R$.
Then we get $I/xA \subset A/xA$. By Lemma \ref{lemma-finite-length-global}
the $R$-module $A/xA$ has finite length as an $R$-module. Hence
$I/xA$ has finite length as an $R$-module. Hence $I$ is finitely
generated as an ideal in $A$.
\end{proof}


\begin{lemma}
\label{lemma-when-finite-module-complete-over-complete-ring}
Let $R$ be a ring. Let $I$ be an ideal of $R$.
Let $M$ be an $R$-module.
If (a) $R$ is $I$-adically complete, (b) $M$ is a finite $R$-module,
and (c) $\bigcap I^nM = (0)$, then $M$ is $I$-adically complete.
\end{lemma}

\begin{proof}
By Lemma \href{https://stacks.math.columbia.edu/tag/0315}{lemma completion generalities (Tag 0315)}
the map $M = M \otimes_R R = M \otimes_R R^\wedge \to M^\wedge$
is surjective. The kernel of this map is $\bigcap I^nM$ hence zero
by assumption. Hence $M \cong M^\wedge$ and $M$ is complete.
\end{proof}


\begin{lemma}
\label{lemma-finite-over-complete-ring}
Let $R$ be a ring. Let $I \subset R$ be an ideal. Let $M$ be an $R$-module.
Assume
\begin{enumerate}
\item $R$ is $I$-adically complete,
\item $\bigcap_{n \geq 1} I^nM = (0)$, and
\item $M/IM$ is a finite $R/I$-module.
\end{enumerate}
Then $M$ is a finite $R$-module.
\end{lemma}

\begin{proof}
Let $x_1, \ldots, x_n \in M$ be elements whose images in $M/IM$ generate
$M/IM$ as a $R/I$-module. Denote $M' \subset M$ the $R$-submodule
generated by $x_1, \ldots, x_n$. By Lemma \href{https://stacks.math.columbia.edu/tag/0315}{lemma completion generalities (Tag 0315)}
the map $(M')^\wedge \to M^\wedge$ is surjective.
Since $\bigcap I^nM = 0$ we see in particular that $\bigcap I^nM' = (0)$.
Hence by Lemma \ref{lemma-when-finite-module-complete-over-complete-ring}
we see that $M'$ is complete, and we conclude that $M' \to M^\wedge$
is surjective. Finally, the kernel of $M \to M^\wedge$ is
zero since it is equal to $\bigcap I^nM = (0)$.
Hence we conclude that $M \cong M' \cong M^\wedge$
is finitely generated.
\end{proof}


\begin{lemma}[Tate]
\label{lemma-tate-japanese}
\begin{reference}
\href{https://stacks.math.columbia.edu/bibliography}{Bibliography: EGA (Theorem 23.1.3)}
\end{reference}
Let $R$ be a ring.
Let $x \in R$.
Assume
\begin{enumerate}
\item $R$ is a normal Noetherian domain,
\item $R/xR$ is a domain and N-2,
\item $R \cong \lim_n R/x^nR$ is complete with respect to $x$.
\end{enumerate}
Then $R$ is N-2.
\end{lemma}

\begin{proof}
We may assume $x \not = 0$ since otherwise the lemma is trivial.
Let $K$ be the fraction field of $R$. If the characteristic of $K$
is zero the lemma follows from (1), see
Lemma \href{https://stacks.math.columbia.edu/tag/032M}{lemma domain char zero N 1 2 (Tag 032M)}. Hence we may assume
that the characteristic of $K$ is $p > 0$, and we may apply
Lemma \ref{lemma-domain-char-p-N-1-2}. Thus given $L/K$
a finite purely inseparable field extension we have to show
that the integral closure $S$ of $R$ in $L$ is finite over $R$.

\medskip\noindent
Let $q$ be a power of $p$ such that $L^q \subset K$.
By enlarging $L$ if necessary we may assume there exists
an element $y \in L$ such that $y^q = x$. Since $R \to S$
induces a homeomorphism of spectra (see Lemma \href{https://stacks.math.columbia.edu/tag/0BRA}{lemma p ring map (Tag 0BRA)})
there is a unique prime ideal $\mathfrak q \subset S$ lying
over the prime ideal $\mathfrak p = xR$. It is clear that
$$
\mathfrak q = \{f \in S \mid f^q \in \mathfrak p\} = yS
$$
since $y^q = x$. Observe that $R_{\mathfrak p}$ is a discrete
valuation ring by Lemma \href{https://stacks.math.columbia.edu/tag/00PD}{lemma characterize dvr (Tag 00PD)}. Then
$S_{\mathfrak q}$ is Noetherian by Krull-Akizuki
(Lemma \ref{lemma-krull-akizuki}). Whereupon we conclude
$S_{\mathfrak q}$ is a discrete valuation ring by
Lemma \href{https://stacks.math.columbia.edu/tag/00PD}{lemma characterize dvr (Tag 00PD)} once again.
By Lemma \ref{lemma-finite-extension-residue-fields-dimension-1} we
see that $\kappa(\mathfrak q)/\kappa(\mathfrak p)$ is
a finite field extension. Hence the integral closure
$S' \subset \kappa(\mathfrak q)$ of $R/xR$ is finite over
$R/xR$ by assumption (2). Since $S/yS \subset S'$ this implies
that $S/yS$ is finite over $R$. Note that $S/y^nS$ has a finite
filtration whose subquotients are the modules
$y^iS/y^{i + 1}S \cong S/yS$. Hence we see that each $S/y^nS$
is finite over $R$. In particular $S/xS$ is finite over $R$.
Also, it is clear that $\bigcap x^nS = (0)$ since an element
in the intersection has $q$th power contained in $\bigcap x^nR = (0)$
(Lemma \href{https://stacks.math.columbia.edu/tag/00IP}{lemma intersect powers ideal module zero (Tag 00IP)}).
Thus we may apply Lemma \ref{lemma-finite-over-complete-ring} to conclude
that $S$ is finite over $R$, and we win.
\end{proof}
\end{document}
